% page 416 of the facsimile (image p-415 of the scan)
% block 154.9 x 246.3 mm ; left margin 8.0 mm
\pgc{-12.700mm}{0.000mm}{
\l{\cel{3}408}
\l{}
\l{\sou{opresar}. (\sou{trans}.)\cel{2}I. Jenar (la pektoro) quale per preso}
\l{\cel{3}qua havas kom efekto, ke onu respiras penoze. - II.}
\l{\cel{3}Koaktar per autoritato tiranala. - III. Igar prostra-}
\l{\cel{3}car korpale. - EFIS.}
\l{}
\l{\sou{optativo}. (\sou{gram.)}\cel{2}Modo di la konjugo, qua expresas la}
\l{\cel{3}deziro. - DEFIS.}
\l{}
\l{\sou{optiko}.\cel{2}Cienco pri la fenomeni di lumo e di vidado. DEFIRS}
\l{}
\l{\sou{optimismo}. I. (\sou{filoz.)}\cel{2}Doktrino segun qua deo kreis nia}
\l{\cel{3}mondo kom la maxim bona ek le posibla. - II. Tendenco}
\l{\cel{3}judikar ke irgo esas la maxim bona ek la posiblaji. DEFIRS}
\l{}
\l{\sou{-or.}\cel{2}(\sou{gram.)}\cel{2}Dezinenco di la infinitivo di futureso.}
\l{}
\l{or.\cel{2}(\sou{konjunciono}). I. En la punto quan atingis la rezono.}
\l{\cel{3}- II. (\sou{logiko)}\cel{2}Konjunciono uzata por venigar la minoro,}
\l{\cel{3}pos la majoro, di silogismo. - FIS.}
\l{}
\l{\sou{oro}.\cel{2}(\sou{kemio}).Metalo flava, tre precoza, tre densa, tre duk-}
\l{\cel{3}tila, tre maleebla, nealterebla da aero nek da aquo, quan}
\l{\cel{3}onu uzas (aloyita kun kelka kupro) por facar moneto-peci}
\l{\cel{3}granda-valora, juveli, e c. -\cel{1}\sou{Simb. kem}. Au.\cel{2}FIS.}
\l{}
\l{\sou{oraklo}. I. (en\cel{1}\sou{la epoki antiqua}). Respondo agita, ye la nomo}
\l{\cel{3}di la deajo, a la persono qua venis konsultar lu en lua}
\l{\cel{3}templo. - II.Dico da deo, enuncita da la profeti, la apos-}
\l{\cel{3}toli, e c. - III. (\sou{metaf.}) Respondo, decido, egardata kom}
\l{\cel{3}neeroriva.- DEFIRS.}
\l{}
\l{\sou{orangutano}. (\sou{zool.)}\cel{2}Granda simio sen-kauda, qua similesas}
\l{\cel{3}pasable la homo de la vidopunto \textquotedbl{}formo extera\textquotedbl{}. - L.}
\l{\cel{3}\sou{simia satyrus}. - DEFIS.}
\l{}
\l{\sou{oranjo.}\cel{2}(\sou{bot.)}\cel{3}Frukto de arbusto di qua la foliaro ne vel-}
\l{\cel{3}kas (L.\cel{1}\sou{citrus aurantium)}\cel{1}sukoza, parfumoza, di qua\cel{2}la}
\l{\cel{3}shelo esas flava oree. - DEFIRS.}
\l{}
\l{\sou{oratoro}. Persono qua, profesione, agas diskursi. - EFIRS.}
\l{}
\l{\sou{oratorio}. (\sou{muziko)}.\cel{2}Dramato religiala, akompanata da mu-}
\l{\cel{3}zikajo, destinita a pleeso sen dekoruri nek kostumi tea-}
\l{\cel{3}trala. - DEFIRS.}
\l{}
\l{\sou{orbito}. I. (\sou{astron.}) Kurvo alonge qua iras, propra-move, ula}
\l{\cel{3}korpi cielala (planeti, kometi, e c.) - II. (\sou{anat.)}\cel{2}Osto-}
\l{\cel{3}kavajo en qua esas la okuli. - EFIRS.}
\l{}
\l{\sou{ordo.}\cel{1}(\sou{eklezio katol.)}\cel{1}Sakramento qua grantas un ek la gradi}
\l{\cel{3}en la hierarkio ekleziala. -}
\l{}
\l{\sou{ordeno}. Singla de la grupi (religiala, +chevalierala, advo-}
\l{\cel{3}katala, medikala) importoza, ye qua konsistas klasifikuro.}
\l{\cel{3}DEFIRS.}
}
